\eventsection{METODOLOGIA}

\eventsubsection{Caracteriza\c{c}\~ao do estudo}

O estudo possui natureza aplicada, abordagem quantitativa e finalidade comparativa. O problema foi modelado em duas tarefas paralelas: regress\~ao para estimar o n\'umero mensal de casos e classifica\c{c}\~ao bin\'aria para detectar surtos. Neste trabalho, considera-se surto a ocorr\^encia de valor mensal acima do percentil 0,9 da s\'erie hist\'orica do respectivo munic\'ipio, desde que o total de casos seja positivo, conforme a Equa\c{c}\~ao \ref{eq:surto}.

\begin{equation}
\label{eq:surto}
\text{surto}_{t} = \mathbb{I}(y_t \geq Q_{0,9}^{(m)} \wedge y_t > 0)
\end{equation}

Em que $y_t$ representa o n\'umero de casos no tempo $t$, $Q_{0,9}^{(m)}$ corresponde ao percentil 0,9 da s\'erie do munic\'ipio $m$ e $\mathbb{I}$ denota a fun\c{c}\~ao indicadora.

\eventsubsection{Dados e engenharia de atributos}

Foram utilizados dados municipais do estado do Amazonas, com vari\'aveis epidemiol\'ogicas e clim\'aticas previamente processadas. O pipeline incluiu ordena\c{c}\~ao temporal por munic\'ipio, gera\c{c}\~ao de \textit{lags}, m\'edias m\'oveis, tend\^encia, componentes sazonais senoidais/cossenoidais e intera\c{c}\~oes entre clima e incid\^encia. Tamb\'em foi aplicado tratamento de valores extremos por \textit{capping} baseado em \textit{z-score} dentro de cada munic\'ipio.

\begin{figure}[H]
\centering
\caption{Arquitetura do pipeline de ingest\~ao, modelagem, avalia\c{c}\~ao e apoio \`a decis\~ao.}
\label{fig:arquitetura}
\includegraphics[width=0.85\textwidth]{./figs/INPUT DATA.png}
\artigofonte{elabora\c{c}\~ao pr\'opria, com base nos experimentos do estudo.}
\end{figure}

A Figura \ref{fig:arquitetura} representa a arquitetura do sistema em quatro camadas encadeadas. A primeira corresponde \`a ingest\~ao e ao pr\'e-processamento dos dados, em que s\~ao integradas s\'eries hist\'oricas municipais de casos de mal\'aria e vari\'aveis clim\'aticas, seguidas por rotinas de limpeza, imputa\c{c}\~ao, normaliza\c{c}\~ao e engenharia de atributos, com destaque para defasagens temporais, m\'edias m\'oveis e rotulagem de surtos a partir do quantil de 90\%. A segunda camada forma o n\'ucleo de modelagem, no qual os dados ingeridos s\~ao submetidos \`a valida\c{c}\~ao cruzada temporal por \textit{TimeSeriesSplit} e \`a busca de hiperpar\^ametros por \textit{GridSearchCV}, alimentando um conjunto comparativo de modelos composto por Random Forest, XGBoost e LightGBM. A terceira camada \'e o m\'odulo de avalia\c{c}\~ao, respons\'avel por medir o desempenho tanto na tarefa de regress\~ao, por meio de MAE, RMSE e W-MAPE, quanto na tarefa de classifica\c{c}\~ao de surtos, por meio de matriz de confus\~ao, AUC-ROC, PR-AUC, curva Precision-Recall e an\'alise de import\^ancia local via SHAP. Por fim, a quarta camada concentra interpretabilidade e apoio \`a decis\~ao, convertendo as sa\'idas anal\'iticas em relat\'orios comparativos, import\^ancia de atributos, comunica\c{c}\~ao a autoridades de sa\'ude, carregamento de previs\~oes e visualiza\c{c}\~ao em painel, de forma a aproximar os resultados do uso operacional em vigil\^ancia epidemiol\'ogica.

\eventsubsection{Modelos e protocolo experimental}

Foram comparados tr\^es algoritmos amplamente utilizados em dados tabulares: Random Forest, XGBoost e LightGBM. O Random Forest foi adotado como refer\^encia da fam\'ilia \textit{bagging}, enquanto XGBoost e LightGBM representaram a fam\'ilia \textit{boosting}. O treinamento considerou ajuste de pesos para a classe minorit\'aria, visando reduzir o impacto do desbalanceamento.

A separa\c{c}\~ao entre treino e teste foi temporal e estrita: os anos anteriores a 2024 foram usados para treinamento, enquanto 2024 em diante compuseram o teste final. No conjunto de teste, foram observadas 291 inst\^ancias, sendo 28 surtos e 263 n\~ao surtos. A busca de hiperpar\^ametros foi realizada por \textit{grid search} com \textit{TimeSeriesSplit} de tr\^es divis\~oes no conjunto de treino.

\begin{table}[H]
\centering
\caption{Resumo da configura\c{c}\~ao experimental.}
\label{tab:configuracao}
\tabelafonte
\begin{tabular}{lll}
\toprule
Componente & Configura\c{c}\~ao & Observa\c{c}\~ao \\
\midrule
Valida\c{c}\~ao & \textit{TimeSeriesSplit} (3 dobras) & Preserva ordem temporal \\
Treino principal & Anos anteriores a 2024 & Ajuste de modelos e limiar \\
Teste principal & Anos de 2024 em diante & Avalia\c{c}\~ao final \\
Classe minorit\'aria & 28 surtos em 291 observa\c{c}\~oes & 9,62\% do teste \\
Limiar de decis\~ao & M\'aximo F1 na valida\c{c}\~ao & Aplicado sem reajuste no teste \\
\bottomrule
\end{tabular}
\artigofonte{elabora\c{c}\~ao pr\'opria.}
\end{table}

Para manter comparabilidade entre as fam\'ilias de modelos e custo computacional vi\'avel, foi utilizado um espa\c{c}o de busca reduzido, suficiente para compara\c{c}\~ao controlada, mas n\~ao exaustivo. Essa decis\~ao segue a l\'ogica do estudo original: privilegiar uma avalia\c{c}\~ao metodologicamente consistente entre algoritmos sobre o mesmo recorte temporal, em vez de maximizar individualmente cada modelo.

Outra decis\~ao importante foi priorizar pondera\c{c}\~ao da classe minorit\'aria durante o treinamento, em vez de reamostragem sint\'etica. Essa escolha buscou preservar a estrutura temporal dos dados, embora estrat\'egias como SMOTE e aprendizado custo-sens\'ivel permane\c{c}am relevantes como refer\^encia metodol\'ogica e compara\c{c}\~ao futura \cite{chawla2002smote,elkan2001foundations}.

\begin{table}[H]
\centering
\caption{Grade resumida de hiperpar\^ametros utilizada na busca.}
\label{tab:grade}
\tabelafonte
\begin{tabular}{lll}
\toprule
Modelo/Tarefa & Hiperpar\^ametro & Valores \\
\midrule
XGBoost & \texttt{n\_estimators} & \{100, 200\} \\
XGBoost & \texttt{learning\_rate} & \{0,01; 0,05\} \\
XGBoost & \texttt{max\_depth} & \{3, 5\} \\
LightGBM & \texttt{n\_estimators} & \{100, 200\} \\
LightGBM & \texttt{learning\_rate} & \{0,01; 0,05\} \\
LightGBM & \texttt{max\_depth} & \{3, 5\} \\
Random Forest & \texttt{n\_estimators} & \{100, 200\} \\
Random Forest & \texttt{max\_depth} & \{5, 10\} \\
\bottomrule
\end{tabular}
\artigofonte{elabora\c{c}\~ao pr\'opria com base no protocolo experimental.}
\end{table}

\eventsubsection{M\'etricas de avalia\c{c}\~ao}

Na regress\~ao, foram consideradas as m\'etricas MAE e RMSE. Na classifica\c{c}\~ao, foram usadas AUC-ROC, AUC-PR, F1 e \textit{Brier Score}. As duas primeiras foram tratadas como m\'etricas principais por serem independentes de limiar e mais adequadas para compara\c{c}\~ao global entre modelos em cen\'arios de baixa preval\^encia. O F1 foi mantido como medida complementar de utilidade operacional, enquanto o \textit{Brier Score} foi utilizado para avaliar qualidade probabil\'istica.

Al\'em das m\'etricas principais no recorte base, o protocolo incorporou an\'alises complementares de calibra\c{c}\~ao, sensibilidade ao quantil de defini\c{c}\~ao de surto, import\^ancia de atributos e robustez temporal anual. Essas extens\~oes permitem interpretar n\~ao apenas qual modelo pontua melhor, mas tamb\'em qu\~ao est\'avel e operacionalmente utiliz\'avel ele permanece quando o regime temporal ou a formula\c{c}\~ao do alvo se altera.
